\subsection{Actual ordinary products and circle traces on every lattice-basis fan}
\begin{statement}{P3K.55}{Proved}
Retain the actual positive selected-cusp source loop of P3F.54/P3F.56,
its physical height, fixed scalar, full line and finite-cap maps. For
every integral unimodular matrix \(P\), the two-triangle fan obtained
from OE.1 by \((x,h)\mapsto(Px,h)\) admits the periodic oriented
ordinary plumbing product constructed below. Its complete period
winding matrix and normal first exponents are
\[
 B_P=(b_{ij})=-6P^{-1}e_1e_1^TP,\qquad
 (l_0,l_1,l_2)=(b_{22}-b_{21},b_{11}-b_{12},b_{22}).
\]
Writing \((\xi,\upsilon)^T=P^{-1}e_2\), define
\[
 w=(\xi,\upsilon,-\xi-\upsilon),\qquad
 d=(-\xi-2\upsilon,-2\xi-\upsilon,\xi-\upsilon).
\]
Its actual local orbit and stable-derivative trace degrees are
\[
 I_i=-3w_i+d_i(l_i+3),\qquad W_i=-3w_i-d_i(l_i+3).
\]
In the positive complex meridians and relative dual cycles, its
actual symmetric variation matrix \(V\) and ambient derivative
\(\delta\) satisfy the integral identities
\[
 V(w_je_j-w_0e_0)=I_je_j-I_0e_0,\qquad
 w_j\delta_j-w_0\delta_0=W_j-W_0\quad(j=1,2).
\]
On each native star the ordinary-product comparison is the explicit
pure-left endpoint \(L^6\). The proof supplies every character,
period, unit, normal and chain map used in these formulas. It leaves
the transverse invariants and the general central return unassigned.
\end{statement}
\begin{proof}
\par

\medskip\noindent\textbf{GF.1. The actual character and period changes}

\par\nobreak

Let \(P\) be an integral two-by-two matrix with determinant \(+1\) or \(-1\), and
write \(R=P^{-1}\). The fan here is the image under \(P\) of OE.1: its two
triangles at each lattice point \(n\) have vertices

\[
 n,\ n+Pe_1,\ n+Pe_2,
 \qquad n+Pe_1,\ n+Pe_2,\ n+P(e_1+e_2).
\]

These are the specified two triangles in each period parallelogram.
The homogeneous lattice map is \((x,h)\) to \((Px,h)\); its determinant is
\(\det P\), and its inverse is integral. On the dense torus use the literal
holomorphic character change

\[
 Z'=Z^{R_{11}}V^{R_{12}},\qquad
 V'=Z^{R_{21}}V^{R_{22}},\qquad q'=q.                 \tag{GF.1}
\]

Its inverse uses the rows of \(P\). Thus it extends to an isomorphism
between the toric charts with the two matched fans; on each matched
unimodular triangle it is the identity in the corresponding three
character coordinates. This is a coordinate description of the actual
chosen fan filling. The original source and line coordinates on its
overlap with the fixed exterior are still obtained by the inverse of
(GF.1), so the source and all finite caps are retained.

The original commuting periods are exactly OE.2:

\[
 h:(Z,V,q)\longmapsto(\kappa qZ,\alpha V,q),\qquad
 B:(Z,V,q)\longmapsto(gZ,\eta qV,q),\qquad\eta=\lambda\beta.
\]

Replace their basis by \(h'_j=h^{P_{1j}}B^{P_{2j}}\), for \(j=1,2\).
This has the inverse integral period change \(R\) and therefore gives
the same actual quotient. In the characters (GF.1) the two periods
are

\[
 h'_1:(Z',V')\longmapsto(qA_{11}Z',A_{21}V'),\qquad
 h'_2:(Z',V')\longmapsto(A_{12}Z',qA_{22}V'),           \tag{GF.2}
\]

where every complete multiplier is retained:

\[
 A_{ij}=(\kappa^{P_{1j}}g^{P_{2j}})^{R_{i1}}
                (\alpha^{P_{1j}}\eta^{P_{2j}})^{R_{i2}}.
                                                               \tag{GF.3}
\]

The power of \(q\) in the \(i\)-th character of \(h'_j\) is \((RP)_{ij}\).
This proves the height factors in (GF.2), including their signs.
The units are nonzero on the same protected collars as the original
ones. At \(q=0\) their parameter winding matrix is exactly

\[
 B_P=(b_{ij})=-6P^{-1}e_1e_1^TP,
                    \qquad\operatorname{tr}B_P=-6.   \tag{GF.4}
\]

There is no omitted unit homotopy in this identity. Write the based
ratio of \(\kappa\) as \(\exp(-6i\theta)\) times \(\exp(f_\kappa)\), and write the
based ratios of \(\alpha,g,\eta\) as \(\exp(f_\alpha)\), \(\exp(f_g)\), \(\exp(f_\eta)\).
The four \(f\)'s are their actual based periodic logarithms from F56.
The periodic logarithm of the ratio \(A_{ij}/A_{ij,*}\) after its displayed
angular factor is

\[
 R_{i1}(P_{1j}f_\kappa+P_{2j}f_g)
       +R_{i2}(P_{1j}f_\alpha+P_{2j}f_\eta).          \tag{GF.5}
\]

This formula supplies every subsequent unit homotopy.

\par

\medskip\noindent\textbf{GF.2. The two vertices and the three full normal loops}

\par\nobreak

In the new character coordinates put \((a,b,c)=T^+_0\) and
\((u,v,w)=T^-_0\). Direct substitution in (GF.2) gives

\[
\begin{array}{c|ccc}
 &u&v&w\\ \hline
 C_0&A_{11}A_{21}bc&A_{21}^{-1}ab&A_{11}^{-1}b^{-1}\\
 C_1&A_{12}A_{22}bc&A_{22}^{-1}c^{-1}&A_{12}^{-1}ac\\
 C_2&a^{-1}&ab&ac.
\end{array}                                                   \tag{GF.6}
\]

Each triple product is \(q\). Keep the first vertex coordinates and
use \((u,A_{22}v,A_{11}w)\) at the second, with the corresponding
initial constant units in the reference fibre. The tangent
coordinates on each curve are then literal reciprocals. The
ordered normal pairs are the same as OP.1, and their normal
transition matrices are \(zD_i\) with

\[
 D_0=\operatorname{diag}(A_{22}/A_{21},A_{11}A_{21}),\quad
 D_1=\operatorname{diag}(A_{11}/A_{12},A_{12}A_{22}),\quad
 D_2=\operatorname{diag}(A_{22},A_{11}).                \tag{GF.7}
\]

For the based loops \(K_i=D_{i,*}^{-1}D_i\), the first diagonal
exponents and the sum of the two exponents are therefore

\[
 l_0=b_{22}-b_{21},\qquad l_1=b_{11}-b_{12},\qquad
 l_2=b_{22},\qquad l_i+\widetilde l_i=-6.              \tag{GF.8}
\]

Every determinant has the same even winding \(-6\), including matrices
\(P\) whose entries have arbitrary signs. We now construct the actual
normal contractions, rather than infer them from that parity.

For an integer \(d\) define

\[
 H_d(s,\theta)=
 (\sin\vartheta(s)e^{id\theta}+j\cos\vartheta(s))
                    (\sin\vartheta(s)+j\cos\vartheta(s))^{-1},
                                                               \tag{GF.9}
\]

where \(\vartheta\) is OP.6's fixed smooth radial function. These maps
are smooth, based, and periodic, and satisfy \(H_d(0,\theta)=1\),
\(H_d(1,\theta)=\exp(id\theta)\). Put

\[
 A=H_{-3},\qquad B_i=H_{l_i+3},\qquad
                  T_i(Q)=A Q B_i,
 \quad Q=n_0+j\overline{n_1}.                         \tag{GF.10}
\]

At \(s=1\) its two complex normal factors are \(\exp(il_i\theta)\)
and \(\exp(i(-6-l_i)\theta)\), respectively: left and right quaternion
phases with exponents \(x,y\) give complex exponents \(x+y,x-y\).
Thus (GF.10) is the required contraction of each angular loop.
Multiplying by the diagonal exponentials from (GF.5), with OP.9's
radial cutoff, gives a contraction of the full \(K_i\). In particular
the original units, norms, and reference values remain present.

At \(P=I\) the right factor on \(C_0\) and \(C_2\) is \(H_3\), while OP.8 used
\(H_{-3}^{-1}\). We retain their exact comparison. Set \(a=\sin\vartheta\),
\(b=\cos\vartheta\) and \(n=ai-bk\), a unit imaginary quaternion. The
expansion of (GF.9) is

\[
 H_d=b^2+a^2\cos(d\theta)
       +ab(1-\cos(d\theta))j+a\sin(d\theta)n.
\]

Conjugation by \(n\) fixes \(n\) and sends \(j\) to \(-j\). Consequently
\(nH_d^{-1}n^{-1}=H_{-d}\). The explicit path

\[
 C_t=\cos(\pi t/2)+n\sin(\pi t/2),\qquad
 C_tH_{-3}^{-1}C_t^{-1},\qquad 0\leq t\leq1,
                                                               \tag{GF.10a}
\]

joins OP's right factor to \(H_3\). It is fixed at \(s=0\), where \(H=1\),
and at \(s=1\), where \(n=i\) commutes with the diagonal endpoint. It is
also fixed at both parameter endpoints. Thus it glues with the
protected vertex collars and the carried unit maps to an actual
based homotopy of the two entire plumbing products. This verifies
the comparison to the initial matrices of K51/K52.

Send \((z,n)\) to \((z,K_i(s(|z|),\theta)n)\) on each tube. The exact
collar identity \(zD_{i,*}K_i=zD_i\) glues this to the two vertex
maps. The normal rescaling \(R_t(z,n)=(z,tn)\), followed by the
cutoff flow of its velocity, carries the original higher normal
terms and their protected vertex collars as in OP.3. The curves'
interiors are disjoint, so these three flows have the stated
disjoint supports. This constructs a periodic oriented ordinary
product of the entire labelled plumbing on this actual fan.

Each separate native star is also explicit. Its second coordinates
are, respectively,

\[
\begin{array}{c|ccc}
 C_0&z'=A_{11}w&n'_0=A_{21}v&n'_1=(A_{11}A_{21})^{-1}u\\
 C_1&z'=A_{22}v&n'_0=A_{12}w&n'_1=(A_{12}A_{22})^{-1}u\\
 C_2&z'=u&n'_0=v&n'_1=w.
\end{array}                                                   \tag{GF.11}
\]

The first normal pairs are \((a,c),(a,b),(b,c)\). Each row gives
exactly \(z'=z^{-1}\), \(n'=zn\). These are actual periodic holomorphic
star products with the complete units shown. The ordinary product
comparison on its normal Hopf boundary is

\[
              q_H\longmapsto (U(z)AU(z)^{-1})q_HB_i.
                                                               \tag{GF.12}
\]

The left factor is precisely OP.20 and has degree \(-6\) on the
Hopf-oriented base sphere times the parameter. The right factor
has only radial and parameter variables and has degree zero.
The explicit cylinder extension and orientation of OR.2 give
the pure-left endpoint \(L^6\). Thus on every separate star the
ordinary relative return is the native relative return composed
with \(L^6\); its right integer is unchanged. This is an actual
boundary-product comparison on the given fan.

\par

\medskip\noindent\textbf{GF.3. Physical weights and exact local orbit maps}

\par\nobreak

The physical delta character in the new coordinates has integral
weights \((\xi,\upsilon)^T=Re_2\). At the first vertex its weights on
\((a,b,c)\) are \((-\xi-\upsilon,\xi,\upsilon)\); at the second they are
\((\xi+\upsilon,-\upsilon,-\xi)\). Define the tangent and normal-difference
weights in the above ordered pairs by

\[
\begin{array}{c|rr}
 &w_i&d_i\\ \hline
 C_0&\xi&-\xi-2\upsilon\\
 C_1&\upsilon&-2\xi-\upsilon\\
 C_2&-\xi-\upsilon&\xi-\upsilon.
\end{array}                                                   \tag{GF.13}
\]

The original physical normal action on \(Q\) is left multiplication
by \(\exp(-iw_i\phi/2)\) and right multiplication by \(\exp(id_i\phi/2)\).
In the ordinary product (GF.10) it becomes exactly

\[
 Q\longmapsto
 (A e^{-iw_i\phi/2}A^{-1})Q
                 (B_i^{-1}e^{id_i\phi/2}B_i).         \tag{GF.14}
\]

The full unit version is obtained by the actual conjugations in
(GF.5), with the same based homotopy. Its source point moves by
\(z\) to \(\exp(iw_i\phi)z\). On the positive radial path \(z=\tan\chi\),
write \(U=\cos\chi+j\sin\chi\) and \(E_i=\exp(iw_i\phi/2)\). The moved
Hopf representative is exactly \(E_i^{-1}UE_i\). At \(q_H=1\) the
normal-sphere orbit trace is consequently

\[
 F_i=E_i^{-1}UE_i A E_i^{-1}A^{-1}U^{-1}
                       B_i^{-1}e^{id_i\phi/2}B_i.     \tag{GF.15}
\]

This includes \(w_i=0\): its left \(A\) factors cancel and the result is
the right-axis map in (GF.14). Every expression closes at \(\phi=2\pi\)
because \(w_i\) and \(d_i\) have the same parity. It is periodic in \(\theta\),
and its radial-end values are \(\theta\)-independent.

For completeness the axis degree in (GF.9) is computed directly.
For \(d=0\) the map is constant. For \(d\) nonzero, the inverse images
of \(-i\) under \(H_diH_d^{-1}\) have \(\vartheta=\pi/4\) and \(\exp(id\theta)=-1\).
There are \(|d|\) of them. At each, \(H_d=j\), and differentiation gives

\[
 \partial_s(H_diH_d^{-1})=4\vartheta'(s)k,\qquad
 \partial_\theta(H_diH_d^{-1})=-d j.
\]

Since \((k,j)\) is the positive outward tangent basis at \(-i\), the
degree is \(-d\). For \(H_d^{-1}iH_d\) the first derivative changes
sign and the second remains \(-dj\), giving degree \(d\).
Thus the axes \(AiA^{-1}\) and \(B_i^{-1}iB_i\) have degrees
\(3\) and \(l_i+3\). Conjugating the first by \(U\) preserves its degree
through \(U_t=\cos(t\chi)+j\sin(t\chi)\); both collapsed ends remain
\(\theta\)-independent. The axis \(UiU^{-1}\) has degree zero.

Factor (GF.15) into \(E_i^{-1}\), \(\exp(iw_i\phi\,\operatorname{Ad}_U(i)/2)\),
\(\exp(-iw_i\phi\,\operatorname{Ad}_{UA}(i)/2)\), and
\(\exp(id_i\phi\,\operatorname{Ad}_{B_i^{-1}}(i)/2)\). Integration of the positive
quaternion volume gives degree \(ad\) for an axis of degree \(d\) and
angle \(a\phi/2\), for every integer \(a\). The Maurer product boundary
forms vanish at the radial ends and at the constant central-sign
angle ends; the parameter faces cancel. Therefore its exact local
normal-meridian degree, in path--parameter--circle orientation, is

\[
                         I_i=-3w_i+d_i(l_i+3).         \tag{GF.16}
\]

For the full tangent derivative, use K52's covering group of
\(GL^+\) and its normalized third form. The local normal block in
(GF.14) has stable degree \(\deg(\mathrm{left})-\deg(\mathrm{right})\). The base rotation
has only the \(\phi\) variable. At \(\phi=2\pi\) its Spin endpoint is
\((-1)^{w_i}\), while the normal Spin endpoint has the same sign;
their product is \(+1\). Thus the full tangent lift is closed even
for odd weights. The same exact boundary-form calculation gives

\[
                         W_i=-3w_i-d_i(l_i+3).         \tag{GF.17}
\]

The ambient/local frame comparison extends through the annulus
times \(D^4\) and has zero third pairing with each capped edge trace.
The original nonzero unit factors and higher derivatives are
carried by the actual homotopies above, so these are degrees of
the full source family, with its original physical circle.

\par

\medskip\noindent\textbf{GF.4. The resulting actual integral relations}

\par\nobreak

Choose generic normal sections at the two ends of each edge and
the fixed vertex-link paths of CT.3. Their three coordinate
moduli are nonzero, and the interpolation of squared moduli
preserves all axis-tube inequalities. The physical rotations,
now with the integer weights in (GF.13), preserve these moduli.
The vertex actions are \(\theta\)-independent, because the vertex
unit maps commute with the diagonal physical circle.

Cap each \(\theta\)-independent endpoint by \(D^2_\theta\) times \(S^1_\phi\),
constant in its \(D^2\) variable. The capped edge trace has class
\(I_im_i\) in the common exterior; the vertex trace has zero third
class. For the graph loops \(\gamma_j=a_j-a_0\) and the relative
two-chains \(x_j\) with those boundaries, the actual positive circle
trace and the complex sphere orientation give

\[
 \Delta x_j=w_jD_j-w_0D_0,\qquad j=1,2.               \tag{GF.18}
\]

This follows on the regular-neighbourhood retract: circle--path
orientation on a curve with weight \(w_i\) has degree \(-w_i\), and
the boundary identity is \(\partial\Delta=-\Delta\partial\). The
integral boundary map identifies the relative three-group with
the complex curve sphere basis, exactly as in CT.3. No claim
that the two vectors in (GF.18) form a primitive sublattice is
needed here.

Let \(V\) be the actual ordinary-product variation matrix on this
fan and \(\delta\) its actual ambient stable-derivative vector. The
two exact prism identities proved in CT.3 and ST.1 now give

\[
 V(w_je_j-w_0e_0)=I_je_j-I_0e_0,\qquad
 w_j\delta_j-w_0\delta_0=W_j-W_0,\qquad j=1,2.        \tag{GF.19}
\]

The matrix \(V\) is symmetric by the integral doubled-exterior
intersection proof of K29, since the boundary meridians generate
\(H_3\) and the return fixes its boundary. Equations (GF.19) are
complete integral identities for the actual constructed product.
They leave the transverse variation/derivative entries to be
computed by the remaining global or flop comparison. In particular
this proof does not assign the missing central return or replace
the needed product-reset map by a scalar-family recurrence.

At \(P=I\), \((l_i)=(0,-6,0)\), \((w_i)=(0,1,-1)\),
\((d_i)=(-2,-1,-1)\); the formulas give \(I=(-6,0,0)\) and \(W=(6,-6,6)\).
Substitution of the completely audited initial matrices of K51/K52
verifies both relations, with every sign retained. This equality
is also executed in the accompanying exact certificate.

The exact certificate is \path{code/part3_root/source_cusp_general_fan_product.py}. It includes symbolic character, period, contraction and orientation identities, and all 1384 integral bases with entries in \([-8,8]\). The written proof establishes the statements for every integral unimodular basis.
\end{proof}
